\documentclass[11pt]{article}
\usepackage[margin=1.05in]{geometry}
\usepackage{amsmath,amssymb,amsthm,mathtools}
\usepackage[colorlinks,linkcolor=blue!50!black,citecolor=blue!50!black]{hyperref}
\usepackage{enumitem}

\theoremstyle{plain}
\newtheorem{thm}{Theorem}[section]
\newtheorem{lem}[thm]{Lemma}
\newtheorem{prop}[thm]{Proposition}
\newtheorem{cor}[thm]{Corollary}
\theoremstyle{definition}
\newtheorem{defn}[thm]{Definition}
\newtheorem{rem}[thm]{Remark}

\newcommand{\At}{\widetilde{A}}
\newcommand{\bA}{\mathbf{A}}
\newcommand{\bB}{\mathbf{B}}
\newcommand{\bI}{\mathbf{I}}
\newcommand{\bM}{\mathbf{M}}
\newcommand{\bR}{\mathbf{R}}
\newcommand{\sm}{\smallsetminus}
\newcommand{\ord}{\operatorname{ord}}
\newcommand{\Pen}{\operatorname{Pen}}
\newcommand{\ES}{\mathcal{E}_S}
\DeclareMathOperator{\WX}{W\!X}

\title{Hypothesis $L^{S}$ is a theorem:\\
sharp triangularity of Motegi's type-$B/C$ weight functions\\
after restriction to a boundary divisor}
\author{Clio}
\date{2026-10-20}

\begin{document}
\maketitle

\begin{abstract}
Let $\At^{X}_{\bA,\bB}$ ($X\in\{B,C\}$) be Motegi's normalised type-$X$
trigonometric weight function for $LG(N,2N)$, indexed by subsets
$\bA,\bB\subseteq[N]$ in the all-cardinality normalisation of
\texttt{arXiv:2610.15344}, and let $S=S_m^{\varepsilon}=\{w_m=\varepsilon\}$ be a
boundary divisor.  I prove that
\[
  \bA\ge\bB\ \text{ and }\ m\notin\bA
  \quad\Longrightarrow\quad
  \At^{X}_{\bA,\bB}\big|_{S}\ne0 ,
\]
for every $N$, both types, every $m\in[N]$ and every finite nonzero
$\varepsilon$ --- with no genericity hypothesis on the remaining variables.  This
is Hypothesis~$L^{S}$ of \texttt{2026-10-19-c2-sa-divisors-full-q}, recorded
there as gap \textbf{G1} and as ``the first thing to attempt next''.  It closes
G1 and makes Lemma \texttt{lem:gnz}, the ``every summand is nonzero'' clause of
Theorem~S and Theorem~S$'$ unconditional on this count.

The proof is Rick's Lemma~L degeneration $w_c=T^{-e_c}$, $T\to0^{+}$, at $q=0$.
My previous note asserted that this degeneration ``cannot be run with $w_m$
pinned at $\varepsilon$''.  That assertion was false, and it was false because I
had paraphrased the lemma without reading it: Rick's Lemma~L \emph{already} pins
two coordinates at finite values while the others degenerate, and all the limit
asks of a pinned coordinate is that it be finite.  I state and prove the general
form --- an arbitrary set of \emph{spectator} coordinates, disjoint from the row
index, pinned at arbitrary finite nonzero values --- for which Rick's two-spectator
lemma and Hypothesis~$L^{S}$ are the cases $|\Sigma|=2$ and $|\Sigma|=1$.  As a
by-product the type-$B$ case of sharp triangularity, which was not in the
literature available to me, comes out on the same line as type $C$.
\end{abstract}

\section{Statement}

\subsection{Setup and the three cited facts}

Fix $N\ge1$ and $X\in\{B,C\}$, with $h=h_X$ ($h_C=2$, $h_B=1$).  Write
$F:=\mathbb{C}(q,w_1,\dots,w_N)$ and $\mathcal{P}:=2^{[N]}$.  For
$\bA=\{a_1<\dots<a_n\}$ and $\bB=\{b_1<\dots<b_r\}$ in $\mathcal{P}$ put
\begin{equation}\label{eq:Adef}
  \At^{X}_{\bA,\bB}
  :=\widetilde{\WX}_{N,n,r}\bigl(w_{\bA}\,\big|\,w\,\big|\,\bB\bigr)
  =\frac{\WX_{N,n,r}(u\,|\,w\,|\,\bB)}{E^{X}_{N,n,r}(u,w)},
  \qquad u:=(w_{a_1},\dots,w_{a_n}),
\end{equation}
Motegi's all-cardinality normalisation, with
$\WX=F(\cdot;\beta,\gamma)$ for $(\beta,\gamma)=(1,0)$ in type $C$ and
$\WX=\bigl(\prod_{k}(1+u_k^{-1})\bigr)^{-1}F(\cdot;1,-1)$ in type $B$, where $F$
is the \emph{Motherfunction} of \texttt{2610.15344} (L473--490):
\begin{align}
  F_{N,n,r}(u|w|\bB;\beta,\gamma)
  &=\frac{\prod_{k=1}^{n}u_k^{-(N+n)}}{(n-r)!}
     \sum_{\sigma\in\mathfrak{S}_n}t_\sigma ,
     \qquad v_j:=u_{\sigma(j)},\label{eq:F}\\
  t_\sigma
  &=\prod_{j=1}^{r}\Bigl[(1-q)(v_j^{2}-1)
      \prod_{k=b_j+1}^{N}(v_j-qw_k)\prod_{k=1}^{b_j-1}(v_j-w_k)\Bigr]
    \notag\\
  &\quad\times\prod_{1\le j<k\le r}\frac{qv_j-v_k}{v_j-v_k}(v_jv_k-1)
    \;\times\prod_{j=r+1}^{n}\ \prod_{k=1}^{N}(v_j-qw_k)
    \notag\\
  &\quad\times\prod_{j=r+1}^{n}\prod_{k=1}^{r}\frac{qv_j-v_k}{v_j-v_k}(v_jv_k-1)
    \;\times\prod_{r+1\le j<k\le n}(v_jv_k-q)
    \;\times\prod_{j=r+1}^{n}(\beta v_j-\gamma).\notag
\end{align}
Slots $j\le r$ are the \textbf{heads} and slots $j>r$ the \textbf{tails}; the
head in slot $j$ is matched with $b_j$.  By the stability relation
(\texttt{eq:typeC-stability-s-ge-r}, L2822--2828) the value with $n=|\bA|$ and no
trailing $1$'s is the stable one, and $\At_{\bA,\bB}=0$ for $|\bA|<|\bB|$
(L2843).

Three facts are quoted from \texttt{2610.15344} (deep-read; the engine of
\S\ref{sec:verif} is my own transcription of the displayed formulas, and every
line of \eqref{eq:F} is checked against Motegi's independent diagonal product
\texttt{normalizationWC} L1078 / \texttt{normalizationWB} L1386):
\begin{description}
\item[(M3) Triangularity (Thm 2.6(i), L603--620; order defined L582--589).]
  $\At_{\bA,\bB}=0$ unless $\bA\ge\bB$, where $\bA\ge\bB$ means $n\ge r$ and
  $a_{n-r+j}\ge b_j$ for $j=1,\dots,r$.
\item[(M5) Integrality (L3302--3309 type $C$, L3722--3729 type $B$).]
  $\At_{\bA,\bB}\in R:=\mathbb{C}(q)[w_1^{\pm},\dots,w_N^{\pm}]$.
\end{description}
By \texttt{lem:order} of \texttt{2026-10-19-c2} the order of (M3) is the
tail-counting order: $\bB\le\bA\iff\bB_{\ge k}\le\bA_{\ge k}$ for all $k\in[N]$,
where $\bB_{\ge k}:=\#\{c\in\bB:c\ge k\}$.

Fix $m\in[N]$ and $\varepsilon$ with $\varepsilon^{h}=1$, let
$S=S_m^\varepsilon:=\{w_m=\varepsilon\}$ with uniformiser $t:=w_m-\varepsilon$
and residue field $k\cong\mathbb{C}(q)(w_j:j\ne m)$.  By (M5),
$\At_{\bA,\bB}\in\mathcal{O}$, so $\ord_S\At_{\bA,\bB}\ge0$ and
$\At_{\bA,\bB}|_{S}$ is the image of $\At_{\bA,\bB}$ in $k$, i.e.\ the result of
substituting $w_m=\varepsilon$ in the Laurent polynomial.

\subsection{The theorem}

\begin{thm}[Hypothesis $L^{S}$]\label{thm:LS}
Let $N\ge1$, $X\in\{B,C\}$, $m\in[N]$ and let $\varepsilon\in\mathbb{C}^{\times}$
be arbitrary.  Let $\bA\subseteq[N]\sm\{m\}$ and $\bB\subseteq[N]$ with
$\bA\ge\bB$.  Then
\[
  \At^{X}_{\bA,\bB}\big|_{w_m=\varepsilon}\ \ne\ 0
  \quad\text{in }\ \mathbb{C}(q)\bigl[w_j^{\pm}:j\ne m\bigr],
  \qquad\text{equivalently}\qquad
  \ord_S\At^{X}_{\bA,\bB}=0 .
\]
\end{thm}

Three things about the quantifiers, each of which is strictly more than
\texttt{hyp:L} asked for.
\begin{enumerate}[label=(\roman*),leftmargin=2.2em]
\item \textbf{$\varepsilon$ is arbitrary.}  The admissibility condition
  $\varepsilon^{h}=1$ plays no role; any finite nonzero value will do.  In
  particular the statement holds in type $B$ at $\varepsilon=-1$, where (BC)
  \emph{fails} (gap G4) --- so $L^{S}$ is not where the two types part company.
\item \textbf{No genericity.}  Unlike (BD), Theorem~\ref{thm:LS} is a statement
  about the generic point of $S$ and needs no condition on $q$ or on
  $(w_j)_{j\ne m}$: it asserts nonvanishing in the residue field $k$ itself.
\item \textbf{Both types, all $N$.}  Nothing is conditional on G3.
\end{enumerate}

Theorem~\ref{thm:LS} is the $|\Sigma|=1$ case of the following, which is also
what makes its relationship to Rick's Lemma~L visible.

\begin{thm}[spectator degeneration]\label{thm:spect}
Let $X\in\{B,C\}$, let $\Sigma\subseteq[N]$, and let $c_s\in\mathbb{C}^{\times}$
be arbitrary for $s\in\Sigma$.  Let $0<e_{c_1}<e_{c_2}<\dots$ be strictly
increasing along $[N]\sm\Sigma=\{c_1<c_2<\dots\}$.  Specialise
\[
  q=0,\qquad w_s=c_s\ (s\in\Sigma),\qquad w_c=T^{-e_c}\ (c\notin\Sigma),
\]
and let $T\to0^{+}$.  Then for every $\bA\subseteq[N]\sm\Sigma$ and every
$\bB\subseteq[N]$,
\[
  \At^{X}_{\bA,\bB}\ \longrightarrow\ [\,\bA\ge\bB\,].
\]
\end{thm}

\begin{rem}[what is Rick's and what is mine]\label{rem:credit}
The mechanism is Rick's Lemma~L
(\texttt{peers/rick/proofs/2026-10-17-q488-converse-order-q1-proofs.pdf},
L123--135), which is Theorem~\ref{thm:spect} in type $C$ with $\Sigma=\{a,b\}$
and $(c_a,c_b)=(x,1/x)$, $x$ generic --- the configuration he needs because he
works on the long-root divisor $H_{ab}=\{w_aw_b=1\}$.  His statement reads, in
his words: ``Let $a,b\notin M$, and let $M\ge\kappa$.  Restrict to $H_{ab}$:
$w_a=x$, $w_b=1/x$, with $x$ fixed and generic.  Take $w_c=t^{-e_c}$ for
$c\notin\{a,b\}$ \dots\ Then $\varphi^{0}_{\kappa}(M)\to1$.''  The parenthesis in
his proof that does the work is ``(For $k\in\{a,b\}$ these are
$1-x^{\pm1}/h_j\to1$, since $h_j\to\infty$.)''

So the lemma was always stated \emph{with} pinned coordinates, and with the
pinned indices excluded from the row index --- which is verbatim my hypothesis
$m\notin\bA$.  \S3 of my own brief asserted the opposite (``that degeneration
cannot be run with $w_m$ pinned at $\varepsilon$, so Hypothesis \texttt{hyp:L} is
not proved here'', \texttt{2026-10-19-c2} L508--511), and the assertion was
manufactured by paraphrase: I recorded the degeneration and dropped the
spectators.  What is mine here is (a) the observation that a spectator may be a
\emph{constant} rather than a free generic parameter, since the limit asks only
that it be finite; (b) the type-$B$ case; (c) the normaliser reduction
(Lemma~\ref{lem:q0}) that transports the $q=0$ statement to $\ord_S$ of the
full-$q$ entry in my normalisation $\At$ rather than his $\varphi^0$; and (d) the
exponent identity of Lemma~\ref{lem:ord}, which replaces the term-by-term
bookkeeping with one line.
\end{rem}

\section{Proof}

\subsection{Step 1: the normalisers disappear at $q=0$}

\begin{lem}\label{lem:q0}
For all $N\ge n\ge r\ge0$ and $X\in\{B,C\}$:
\begin{enumerate}[label=(\arabic*),leftmargin=2.2em]
\item $E^{X}_{N,n,r}(u,w)\big|_{q=0}=1$;
\item $\At^{X}_{\bA,\bB}=N/(D_0D_1)$, where $N\in\mathbb{C}[q][w^{\pm}]$, where
  $D_0$ is a nonzero element of
  $\mathbb{C}\bigl[(w_a)^{\pm}_{a\in\bA}\bigr]$ --- involving \emph{no} $q$ and
  \emph{no} $w_c$ with $c\notin\bA$ --- and where $D_1$ is a product of factors
  $(1-q\cdot\text{monomial})$, so that $D_1|_{q=0}=1$;
\item consequently $\At^{X}_{\bA,\bB}$ is regular at $q=0$ and
  $\At^{X}_{\bA,\bB}\big|_{q=0}=\WX_{N,n,r}(u|w|\bB)\big|_{q=0}$;
\item if $m\notin\bA$, $\varepsilon\ne0$, and
  $\bigl(\At^{X}_{\bA,\bB}\big|_{q=0}\bigr)\big|_{w_m=\varepsilon}\ne0$, then
  $\ord_S\At^{X}_{\bA,\bB}=0$.
\end{enumerate}
\end{lem}

\begin{proof}
(1) By definition (\texttt{eq:typeCnormalizationfactor} L1057--1062,
\texttt{eq:typeBnormalizationfactor} L1361--1367) every factor of $E^{X}$ is of
the form $(1-q\cdot\text{monomial})^{\pm1}$, together with the single factor
$(1-q)^{-(n-r)}$.  Each becomes $1$ at $q=0$.

(2) In \eqref{eq:F} every factor of $t_\sigma$ is a polynomial in $q$ and in the
$w$'s except for the denominators $(v_j-v_k)$ of lines 2a and 3a; these, and the
prefactor $\prod_k u_k^{-(N+n)}$, involve only the $u_j=w_{a_j}$ with
$a_j\in\bA$, and no $q$.  Putting the $n!$ terms over the common denominator
$\prod_{j<k}(u_j-u_k)\cdot\prod_k u_k^{\,N+n}$ gives
$F=P/\bigl(\prod_{j<k}(u_j-u_k)\prod_k u_k^{N+n}\bigr)$ with
$P\in\mathbb{C}[q][w^{\pm}]$; in type $B$ there is the further $q$-free factor
$\prod_k(1+u_k^{-1})^{-1}=\prod_k u_k/(u_k+1)$, again in the $u$'s only.  Finally
$1/E^{X}$ is, by (1)'s description of $E^{X}$, a quotient of two products of
factors $(1-q\cdot\text{monomial})$ and powers of $(1-q)$; absorb its numerator
into $N$ and its denominator into $D_1$.  Set $D_0:=\prod_{j<k}(u_j-u_k)\prod_k
u_k^{N+n}$ (times $\prod_k(u_k+1)$ in type $B$), which is nonzero and lies in
$\mathbb{C}[(w_a)^{\pm}_{a\in\bA}]$.

(3) By (2), $D_0D_1|_{q=0}=D_0\ne0$, so $\At^{X}_{\bA,\bB}$ is regular at $q=0$;
and $1/E^{X}|_{q=0}=1$ by (1).

(4) We prove the contrapositive.  Suppose $\ord_S\At_{\bA,\bB}\ge1$, i.e.\
$\At_{\bA,\bB}$ maps to $0$ in the residue field $k$ of the local ring at the
generic point of $S$.  Write $\At_{\bA,\bB}=N/(D_0D_1)$ as in (2) and substitute
$w_m=\varepsilon$.  Since $m\notin\bA$, $D_0$ does not involve $w_m$ and
$D_0|_{w_m=\varepsilon}=D_0\ne0$; and $D_1|_{w_m=\varepsilon}$ is a product of
factors $(1-q\cdot M)$ with $M$ a Laurent monomial in $(w_j)_{j\ne m}$ times a
power of $\varepsilon\ne0$, hence a nonzero element of
$\mathbb{C}[q][w_j^{\pm}:j\ne m]$.  So the vanishing of $\At_{\bA,\bB}$ in $k$
forces $N|_{w_m=\varepsilon}=0$ in $\mathbb{C}[q][w_j^{\pm}:j\ne m]$.  A
polynomial identity survives any further substitution, so
$N|_{w_m=\varepsilon}\big|_{q=0}=0$ as well.  Therefore, using
$D_1|_{q=0}=1$ and $D_0|_{q=0,\,w_m=\varepsilon}=D_0\ne0$,
\[
  \bigl(\At_{\bA,\bB}\big|_{q=0}\bigr)\big|_{w_m=\varepsilon}
  =\frac{N\big|_{q=0,\,w_m=\varepsilon}}{D_0}
  =0 .
\]
(The two substitutions commute because neither $D_0$ nor $D_1$ vanishes under
either one.)
\end{proof}

So it suffices to prove Theorem~\ref{thm:spect}: it exhibits a point --- a point
at the end of a curve, which is all a nonvanishing statement needs --- at which
$\At_{\bA,\bB}|_{q=0}|_{w_m=\varepsilon}$ equals $1$, hence is nonzero, and
Lemma~\ref{lem:q0}(4) converts that into $\ord_S\At_{\bA,\bB}=0$.
Note that $\varepsilon$ enters only as the pinned value, so $\varepsilon^{h}=1$
is never used.

\subsection{Step 2: the specialised terms}

Fix the data of Theorem~\ref{thm:spect}.  Extend $e$ to all of $[N]$ by setting
\begin{equation}\label{eq:E}
  E(c):=e_c\ \ (c\notin\Sigma),\qquad E(s):=0\ \ (s\in\Sigma),
\end{equation}
so that $w_c\asymp T^{-E(c)}$ for every $c$ ($w_s=c_s$ is $T^{0}$ exactly).  Thus
$E>0$ off $\Sigma$, $E=0$ on $\Sigma$, and $E$ is \emph{strictly increasing} on
$[N]\sm\Sigma$.  Since $\bA\cap\Sigma=\emptyset$:
\begin{equation}\label{eq:distinct}
  E\ \text{is injective on }\bA,\qquad
  E(c)>E(k)\ \text{ whenever } c\in\bA,\ k<c .
\end{equation}
The second clause holds because either $k\notin\Sigma$, and then $E(k)=e_k<e_c=E(c)$
by strict monotonicity, or $k\in\Sigma$, and then $E(k)=0<E(c)$.  \emph{This one
line is where the spectators are absorbed}, and it is the only place they are
mentioned again.

At $q=0$ the Motherfunction simplifies: $\prod_{k>b_j}(v_j-qw_k)=v_j^{N-b_j}$,
$\prod_{k=1}^{N}(v_j-qw_k)=v_j^{N}$, $(qv_j-v_k)=-v_k$ and
$(v_jv_k-q)=v_jv_k$, so
\begin{align}
  t_\sigma\big|_{q=0}
  &=\prod_{j=1}^{r}\Bigl[(v_j^{2}-1)\,v_j^{\,N-b_j}
        \prod_{k=1}^{b_j-1}(v_j-w_k)\Bigr]
   \ \times\!\!\prod_{1\le j<k\le r}\!\frac{-v_k(v_jv_k-1)}{v_j-v_k}
   \label{eq:t0}\\
  &\quad\times\prod_{j=r+1}^{n}v_j^{\,N}\,(\beta v_j-\gamma)
   \ \times\!\!\prod_{j=r+1}^{n}\prod_{k=1}^{r}\!\frac{-v_k(v_jv_k-1)}{v_j-v_k}
   \ \times\!\!\prod_{r+1\le j<k\le n}\!v_jv_k .\notag
\end{align}
Observe already that \textbf{the only surviving dependence on $w_c$ for
$c\notin\bA$ is through $\prod_{j\le r}\prod_{k<b_j}(v_j-w_k)$}: the two places
where a general inhomogeneity could enter carried a factor $q$ and have become
pure powers of $v_j$.

Write $\sigma$ as the data of an ordered head tuple
$(\eta_1,\dots,\eta_r)\in\bA^{r}$ of distinct elements, $v_j=w_{\eta_j}$, together
with the remaining $n-r$ elements $\tau\subseteq\bA$ as tails in some order.  The
tail-dependent factors of \eqref{eq:t0} are symmetric in the tails, so the
$(n-r)!$ terms with the same $(\eta,\tau)$ are equal and are divided out by the
$(n-r)!$ in \eqref{eq:F}.

\begin{lem}[vanishing terms]\label{lem:zero}
$t_\sigma|_{q=0}\equiv0$ if and only if $\eta_j<b_j$ for some $j\le r$.
\end{lem}

\begin{proof}
If $\eta_j<b_j$ then $k=\eta_j$ occurs in $\prod_{k<b_j}(v_j-w_k)$, giving the
factor $w_{\eta_j}-w_{\eta_j}=0$.  Conversely suppose $\eta_j\ge b_j$ for all
$j$.  Then no factor of \eqref{eq:t0} vanishes under the specialisation: the
$w_c$ with $c\notin\Sigma$ are distinct powers of $T$ and the $c_s$ are nonzero,
so $v_j\ne0$, $v_j^{2}\ne1$ and $v_jv_k\ne1$ for $T$ small, while
$v_j-w_k\ne0$ for $k<b_j\le\eta_j$ by \eqref{eq:distinct} ($E(v_j)>E(k)$ forces
$|v_j|>|w_k|$ for $T$ small, and in particular $v_j\ne w_k$).  Also
$\beta v_j-\gamma\ne0$ in both types since $|v_j|\to\infty$.
\end{proof}

\subsection{Step 3: one exponent identity}

For a term $\sigma$ with heads $(\eta_1,\dots,\eta_r)$ and tail set $\tau$ define
the \textbf{penalty}
\begin{equation}\label{eq:pen}
  \Pen(\sigma):=\sum_{1\le j<k\le r}\bigl(E(\eta_j)-E(\eta_k)\bigr)^{+}
   \;+\;\sum_{c\in\tau}\ \sum_{k=1}^{r}\bigl(E(c)-E(\eta_k)\bigr)^{+}
   \ \ \ge0 ,
\end{equation}
where $x^{+}:=\max(x,0)$.  Thus $\Pen$ charges each inverted head pair and each
tail that sits above a head.

\begin{lem}[exponent identity]\label{lem:ord}
Let $\sigma$ be a term with $t_\sigma|_{q=0}\not\equiv0$, i.e.\ $\eta_j\ge b_j$
for all $j$.  Under the specialisation of Theorem~\ref{thm:spect},
\[
  \ord_T\Bigl(\textstyle\prod_{k=1}^{n}u_k^{-(N+n)}\cdot t_\sigma|_{q=0}\Bigr)
  \;=\;\Pen(\sigma),
\]
and if $\Pen(\sigma)=0$ the leading coefficient is $+1$.
\end{lem}

\begin{proof}
Each factor of \eqref{eq:t0} is a binomial $Y-Z$ (or a ratio of two such) in which
one monomial strictly dominates the other as $T\to0^{+}$, so
$\ord_T(Y-Z)=\min(\ord_T Y,\ord_T Z)$ and the leading coefficient is that of the
dominant monomial.  Write $S_H:=\sum_{j\le r}E(\eta_j)$,
$S_\tau:=\sum_{c\in\tau}E(c)$ and $S_{\bA}:=S_H+S_\tau=\sum_{c\in\bA}E(c)$.
Recall $\ord_T w_c=-E(c)$ and that $|v_j|\to\infty$ for every slot, since
$\bA\cap\Sigma=\emptyset$.

\emph{Heads.}  Slot $j\le r$ contributes
$(v_j^{2}-1)\sim v_j^{2}$, then $v_j^{N-b_j}$ exactly, then
$\prod_{k<b_j}(v_j-w_k)$.  By \eqref{eq:distinct} and $k<b_j\le\eta_j$ we have
$E(\eta_j)>E(k)$, so \emph{each} of these $b_j-1$ factors is $\sim v_j$, with
leading coefficient $1$.  (This is the step that consumes the spectators: for
$k\in\Sigma$ the factor is $v_j-c_k\sim v_j$ because $|v_j|\to\infty$.)  Total:
\[
  \ord_T=-\bigl(2+(N-b_j)+(b_j-1)\bigr)E(\eta_j)=-(N+1)E(\eta_j),
\]
leading coefficient $1$, and --- note --- \emph{independent of $b_j$}.

\emph{Tails.}  Slot $j>r$ contributes $v_j^{N}$, then $(\beta v_j-\gamma)\sim\beta
v_j=v_j$ (type $C$: $\gamma=0$, exactly; type $B$: $v_j+1\sim v_j$), and appears
in $n-r-1$ of the factors $v_jv_k$ of the last product.  Total
$\ord_T=-(N+1+n-r-1)E(c)=-(N+n-r)E(c)$, leading coefficient $1$.

\emph{Pairs.}  Both pair products have the shape
$\dfrac{-v_k(v_jv_k-1)}{v_j-v_k}$ with the \emph{second} slot $k$ a head.  Put
$\alpha:=E(v_j)$, $\beta':=E(v_k)$.  The numerator is $\sim -v_jv_k^{2}$ with
$\ord_T=-(\alpha+2\beta')$; the denominator is $\sim-v_k$ if $\alpha<\beta'$ and
$\sim v_j$ if $\alpha>\beta'$, so $\ord_T=-\max(\alpha,\beta')$.  Hence
\[
  \ord_T\Bigl(\tfrac{-v_k(v_jv_k-1)}{v_j-v_k}\Bigr)
  =-(\alpha+2\beta')+\max(\alpha,\beta')
  =-(\alpha+\beta')+(\alpha-\beta')^{+},
\]
and when $\alpha<\beta'$ the ratio of leading monomials is
$(-v_jv_k^{2})/(-v_k)=v_jv_k$, leading coefficient $+1$.  By
\eqref{eq:distinct} $\alpha\ne\beta'$ always, so no cancellation of leading
monomials can occur.

\emph{Summation.}  Splitting off the $(\cdot)^{+}$ terms, which assemble to
exactly \eqref{eq:pen}:
\begin{align*}
  \ord_T t_\sigma\big|_{q=0}
  &= -(N+1)S_H-(N+n-r)S_\tau
     -\!\!\sum_{j<k\le r}\!\!\bigl(E(\eta_j)+E(\eta_k)\bigr)
     -\!\!\sum_{c\in\tau}\sum_{k\le r}\!\bigl(E(c)+E(\eta_k)\bigr)+\Pen(\sigma)\\
  &= -(N+1)S_H-(N+n-r)S_\tau-(r-1)S_H-\bigl(rS_\tau+(n-r)S_H\bigr)+\Pen(\sigma)\\
  &= -\bigl[(N+1)+(r-1)+(n-r)\bigr]S_H-\bigl[(N+n-r)+r\bigr]S_\tau+\Pen(\sigma)\\
  &= -(N+n)S_H-(N+n)S_\tau+\Pen(\sigma)\;=\;-(N+n)S_{\bA}+\Pen(\sigma).
\end{align*}
The prefactor has $\ord_T\prod_k u_k^{-(N+n)}=+(N+n)S_{\bA}$, which cancels the
first term --- note $S_{\bA}$ does not depend on $\sigma$.  The leading
coefficient is the product of the leading coefficients listed above, all $+1$,
provided every pair is in its $\alpha<\beta'$ case, which is exactly
$\Pen(\sigma)=0$.  In type $B$ there is in addition the prefactor
$\prod_k(1+u_k^{-1})^{-1}\to1$.
\end{proof}

\begin{lem}[the dominant term]\label{lem:dom}
$\Pen(\sigma)=0$ if and only if $(\eta_1,\dots,\eta_r)$ is the increasing
enumeration of the $r$ largest elements of $\bA$ and $\tau$ is the set of the
$n-r$ smallest.  Such a $\sigma$ exists and is unique up to the ordering of
$\tau$, and the corresponding term is nonzero if and only if $\bA\ge\bB$.
\end{lem}

\begin{proof}
By \eqref{eq:distinct}, $E$ is injective and increasing on $\bA$, so
$\Pen(\sigma)=0$ forces $E(\eta_1)<\dots<E(\eta_r)$, i.e.\ $\eta_1<\dots<\eta_r$,
and $E(c)<E(\eta_k)$ for every tail $c$ and every head, i.e.\ every tail is
smaller than every head.  Those two conditions say precisely that the head set is
the top $r$ of $\bA$, listed increasingly.  Conversely that choice gives
$\Pen=0$.  For such $\sigma$, $\eta_j=a_{n-r+j}$, so by Lemma~\ref{lem:zero} the
term is nonzero iff $a_{n-r+j}\ge b_j$ for all $j$, which is (M3)'s order
$\bA\ge\bB$.
\end{proof}

\subsection{Step 4: conclusion}

\begin{proof}[Proof of Theorem~\ref{thm:spect}]
By Lemma~\ref{lem:q0}(3), $\At^{X}_{\bA,\bB}|_{q=0}=\WX|_{q=0}$, which by
\eqref{eq:F} is $\frac{1}{(n-r)!}\prod_k u_k^{-(N+n)}\sum_\sigma t_\sigma|_{q=0}$
(times $\prod_k(1+u_k^{-1})^{-1}\to1$ in type $B$).  By Lemma~\ref{lem:zero} only
terms with $\eta_j\ge b_j$ for all $j$ contribute, and for those
Lemma~\ref{lem:ord} gives $\ord_T=\Pen(\sigma)\ge0$.

If $\bA\ge\bB$: by Lemma~\ref{lem:dom} there are exactly $(n-r)!$ terms with
$\Pen=0$ --- the $(n-r)!$ tail orderings of one head choice --- each nonzero with
leading coefficient $+1$ and $\ord_T=0$; after dividing by $(n-r)!$ they
contribute $1+O(T^{\delta})$.  Every other nonzero term has $\Pen(\sigma)>0$ and
hence $\ord_T>0$, so contributes $O(T^{\delta})$ for some $\delta>0$.  Therefore
$\At^{X}_{\bA,\bB}\to1$.

If $\bA\not\ge\bB$: the $\Pen=0$ terms vanish identically by
Lemma~\ref{lem:dom}, every surviving term has $\ord_T>0$, and
$\At^{X}_{\bA,\bB}\to0$.  (For $|\bA|<|\bB|$ the entry is $0$ outright, L2843.)
\end{proof}

\begin{proof}[Proof of Theorem~\ref{thm:LS}]
Take $\Sigma=\{m\}$, $c_m=\varepsilon$, and any strictly increasing
$(e_c)_{c\ne m}$.  Since $\bA\ge\bB$, Theorem~\ref{thm:spect} gives
$\At^{X}_{\bA,\bB}|_{q=0}\to1$ along the curve
$\{w_m=\varepsilon,\ w_c=T^{-e_c}\}$.  A rational function that is identically
zero on $\{w_m=\varepsilon\}$ is zero along that curve, where it is defined
(Lemma~\ref{lem:q0}(2): the denominator involves only $(w_a)_{a\in\bA}$, which are
distinct powers of $T$); its limit would be $0$, not $1$.  Hence
$\At^{X}_{\bA,\bB}|_{q=0}|_{w_m=\varepsilon}\ne0$, and
Lemma~\ref{lem:q0}(4) gives $\ord_S\At^{X}_{\bA,\bB}=0$.
\end{proof}

\begin{cor}[sharp triangularity, both types, all $N$]\label{cor:sharp}
For all $\bA,\bB\subseteq[N]$ and $X\in\{B,C\}$:
$\At^{X}_{\bA,\bB}\ne0$ in $F$ if and only if $\bA\ge\bB$.
\end{cor}
\begin{proof}
$\Sigma=\emptyset$ in Theorem~\ref{thm:spect} (``$\Leftarrow$''); ``$\Rightarrow$''
is (M3).
\end{proof}

\begin{rem}[why $m\notin\bA$ is not removable, and why $m\in\bB$ is harmless]
\label{rem:hyp}
The hypothesis $m\notin\bA$ is used exactly once, in \eqref{eq:distinct}: it makes
$E$ injective on $\bA$ and makes every head blow up.  If $m\in\bA$ then $w_m$ is
one of the $v_j$, it stays finite while its neighbours blow up, $E$ is no longer
injective on $\bA$, and the pair analysis loses its strict comparison.  That is
not an artefact: (BD) says $\ord_S\At_{\bI,\bI}=1$ when $m\in\bI$, so the
conclusion is false for $\bA=\bB\ni m$.  \textbf{The hypothesis belongs to the
step that separates the two cases, which is the injectivity of $E$ on the row
index.}

By contrast $m\in\bB$ --- the column meeting the pinned coordinate, which the
brief's gate 0b isolated as its own sub-family and which I expected to split ---
is \emph{invisible} to the proof: $\bB$ enters only through the integers $b_j$ in
the head exponents, and Lemma~\ref{lem:ord} showed the head contribution
$-(N+1)E(\eta_j)$ is independent of $b_j$.  The slot $j$ with $b_j=m$ simply has
one fewer factor in $\prod_{k<b_j}$ and one more in $v_j^{N-b_j}$.  So the four
cells of gate 0b have one answer after all, by a mechanism that explains why; the
cell structure is real combinatorially (\S\ref{sec:verif}) but carries no
analytic content.
\end{rem}

\section{Consequences for the S-divisor program}
\label{sec:cons}

With Theorem~\ref{thm:LS} in hand, in \texttt{2026-10-19-c2-sa-divisors-full-q}:
\begin{itemize}[leftmargin=1.4em]
\item \textbf{G1 closes.}  Hypothesis \texttt{hyp:L} becomes a theorem, in both
  types, with no genericity and no restriction on $\varepsilon$.
\item \texttt{lem:gnz} ($\gamma^{\nu}_{\bM}\ne0$ for $\bM\le\nu$, $m\in\bM$) now
  rests only on (BD) and \texttt{lem:inverse}.  It applies $L^{S}$ at
  $\bA=\bM^{\mathrm c}$, $\bB=\nu^{\mathrm c}$ --- a restricted instance; the
  general theorem covers it, and Remark~\ref{rem:hyp} explains why the restricted
  instance was never easier.
\item The ``every summand is nonzero'' clause of \texttt{thm:S}, and
  \texttt{thm:Sprime} (the boundary converse at full $q$), lose their dependence
  on G1.  In type $C$ they are now unconditional modulo \texttt{lem:inverse}
  (G5, Motegi's biorthogonality).  In type $B$ they remain conditional on
  \textbf{G3} --- (BD) in type $B$ --- which is an independent gap:
  Theorem~\ref{thm:LS} does \emph{not} supply it, because (BD) is the statement
  $\ord_S\At_{\bI,\bI}=1$ for $m\in\bI$, i.e.\ exactly the case $m\in\bA$ that
  Remark~\ref{rem:hyp} excludes.
\item \textbf{The two open gaps are now separated.}  The falsifier recorded at
  \texttt{rem} L941ff of that paper said: a false negative at fixed $\chi$ that
  survives a change of $\chi$ refutes $L^{S}$, while one that moves with $\chi$ is
  G2 (cancellation).  Since $L^{S}$ is true, \emph{every} false negative is now
  G2.  The discriminating experiment of \S7 of the brief has become a
  one-sided measurement of G2 alone.
\end{itemize}
Corollary~\ref{cor:sharp} also discharges, in type $B$ and from my side, the
sharp-triangularity input that Rick's Lemma~L supplies in type $C$.

\section{Computational verification}
\label{sec:verif}

All computations are pure Python + SymPy (SageMath is not installed in this
container), exact rational arithmetic throughout, in
\texttt{proofs/code-1020-LS/}.  The engine \texttt{q484\_engine.py}
(\texttt{md5 f53c0258\dots}) is my own transcription of \eqref{eq:F} and of
Motegi's normalisers, cross-checked against his independent positive-root product
\texttt{diag\_rootproduct}.  Every domain is stated before its number.

\paragraph{(V0) The hypothesis class, split before measuring
(\texttt{cells.py}).}  Exhaustive over $N\le5$, every $m$, every pair
$(\bA,\bB)$ with $\bA\ge\bB$ and $m\notin\bA$.  Per type-$\varepsilon$ instance at
$N\le4$: $128$ in cell $m\notin\bB,\bA\ne\bB$; $56$ in $m\in\bB,\bA\ne\bB$; $49$
diagonal.  Total $233$, and $233\times3=\mathbf{699}$, $393\times3=\mathbf{1179}$
--- the census and plant sizes of the previous paper, reproduced exactly from
combinatorics alone, which confirms that the recorded domain is the intended one.
No cell is empty for $N\ge2$.  One combinatorial fact fell out and is used
nowhere but is worth recording: if $m=N$ the cell $m\in\bB$ is \emph{empty}, since
$N\in\bB$ forces $\bA_{\ge N}\ge1$, i.e.\ $N\in\bA$.

\paragraph{(V1) The degeneration limit (\texttt{degen.py}).}
Prediction written before the run: the limit is $1$ on $\bA\ge\bB$ and $0$
otherwise --- the \emph{value}, not merely nonvanishing.
Domain: $N=1,\dots,5$; $X\in\{B,C\}$; every $m\in[N]$; $\varepsilon=\pm1$ (type
$C$) and $+1$ (type $B$); \emph{every} pair $(\bA,\bB)$ with $m\notin\bA$,
exhaustive; $T=10^{-1},10^{-2},10^{-3}$, exact rationals.  Convergence is graded
by the decay exponent $p$ in $|v-\text{pred}|\asymp cT^{p}$, requiring $p>1/2$,
because the true decay is $O(T)$ and a fixed $10^{-3}$ threshold manufactured $11$
false failures on the first run.
\[
  \textbf{9558/9558:}\quad
  3321\ \text{instances with }\bA\ge\bB\ \to1,\qquad
  6237\ \text{plant instances}\ \to0 .
\]
At $N\le4$ these are $699$ and $1179$: \emph{the same two numbers as the previous
paper's full-$q$ census, reached by a different mechanism} ($q=0$ degeneration
limit versus $\ord_t$ of the full-$q$ entry).

\paragraph{(V2) Plant on the hierarchy (\texttt{degen.py --plant-ord}).}
The same computation with $e$ \emph{reversed}, so that
$w_1\gg\dots\gg w_N$: $12$ instances deviate and their errors \emph{grow} with
decreasing $T$.  The monotonicity of $e$ is therefore load-bearing and the
instrument can see its absence --- without this control the census would not
distinguish ``the hierarchy produces the limit'' from ``the limit is $1$ anyway''.

\paragraph{(V3) The exponent identity, term by term (\texttt{terms.py}).}
This dissects what (V1) sums, so it is an independent reading of
Lemma~\ref{lem:ord} rather than a second reading of Theorem~\ref{thm:spect}.
Domain: $N=2,3,4$, both types, every $m$, every $\varepsilon$, every
$(\bA,\bB)$ with $\bA\ge\bB$ and $m\notin\bA$, \emph{every} $\sigma\in
\mathfrak{S}_n$; $\ord_T$ and leading coefficient computed exactly as a Laurent
rational function in a symbolic $T$.
\begin{itemize}[leftmargin=1.4em]
\item $1119/1119$ nonzero terms satisfy $\ord_T=\Pen(\sigma)$;
\item $612$ identically-zero terms, \emph{all} predicted by
  Lemma~\ref{lem:zero}'s criterion ``$\eta_j<b_j$ for some $j$'';
\item $819$ terms have $\Pen=0$, all with leading coefficient exactly $+1$, and
  the number of $\Pen=0$ terms per instance is exactly $(n-r)!$ in every one of
  the instances --- which is Lemma~\ref{lem:dom}'s uniqueness claim, measured.
\end{itemize}
A first version of this check also demanded leading coefficient $+1$ for
\emph{sub}dominant terms.  That is not claimed and is false: $243$ of them have
coefficient $-1$.  Recorded because the orders matched in all $1119$ cases on that
run too, so the only thing the ``failure'' measured was my over-statement.

\paragraph{(V4) The general spectator form (\texttt{spect.py}).}
Domain: $N\le4$, both types, \emph{every} $\Sigma\subseteq[N]$ (including
$\emptyset$ and $|\Sigma|\ge2$), every $(\bA,\bB)$ with $\bA\cap\Sigma=\emptyset$,
spectator values drawn from the deliberately non-generic list
$\{1,-1,5/3,3/5\}$ --- which includes both roots of unity and a reciprocal pair,
i.e.\ Rick's configuration.  $1068$ limit-$1$ and $2040$ plant instances:
$\mathbf{3108/3108}$.

\paragraph{(V5) The normaliser lemma, Lemma~\ref{lem:q0} (\texttt{denom.py}).}
$E^{X}_{N,n,r}|_{q=0}=1$ verified symbolically for all $N\le4$, all $n\ge r\ge0$,
both types: all cases, no exceptions.  The load-bearing half of
Lemma~\ref{lem:q0}(2) --- that the denominator of $\At^{X}_{\bA,\bB}$ in lowest
terms, evaluated at $q=0$, is nonzero and involves no $w_c$ with $c\notin\bA$ ---
verified exhaustively at $N\le3$, both types, every $(\bA,\bB)$ with
$\bA\ge\bB$: $\mathbf{96/96}$.  In all $96$ it is in fact a single monomial,
which is stronger than the lemma claims; the lemma is deliberately stated in the
weaker form that the proof consumes, because the stronger one depends on
cancellations among the $n!$ terms that I have not proved.

\paragraph{(V6) Full-$q$ direct census at $N=5$ (\texttt{fullq\_N5.py}).}
Independent of (V1)--(V4): it reads $\ord_t\At_{\bA,\bB}$ at $w_m=\varepsilon+t$
with $q=3/11$ and the other $w_j$ at a rational point \emph{gated} on (BD) ---
require $\ord_t$ of Motegi's independent diagonal root product to equal
$[m\in\bI]$ for \emph{every} $\bI$, and redraw otherwise.  That gate is the one
whose absence produced a confident, legible and wholly false
$\ord_S\det\Phi=8$ on 2026-10-19, when the drawn point happened to satisfy
$w_3=q^{-1}$.  Domain: $N=5$, both types, every $m\in[5]$, every admissible
$\varepsilon$, exhaustive over $(\bA,\bB)$ with $m\notin\bA$.  Prediction written
first: $\ord_t=0$ on $\bA\ge\bB$, identically zero otherwise.
\[
  \textbf{2622/2622}\ \ (\ord_t=0),\qquad
  \textbf{5058/5058}\ \ \text{plant}.
\]
The two totals are forced by (V0): $\sum_{m=1}^{5}|\{\bA\ge\bB,\,m\notin\bA\}|
=210+196+181+161+126=874$ per $(X,\varepsilon)$ arm and $874\times3=2622$;
likewise $1686\times3=5058$.  They came out equal, so the census covered the
domain it reports and not a truncation of it.  This extends the measured domain
of $L^{S}$ from $N\le4$ to $N=5$.

\paragraph{Summary of the verification.}  Six instruments, three of them
independent mechanisms for the same claim: the $q=0$ degeneration limit (V1,
$9558/9558$ at $N\le5$), the term-by-term exponent dissection of what that limit
sums (V3, $1119/1119$ plus $612/612$ predicted zeros), and the full-$q$
$\ord_t$ reading that never mentions $q=0$ at all (V6, $2622/2622$ at $N=5$).
Two plants fire (V2 on the hierarchy, and the $\bA\not\ge\bB$ plant in V1 and
V6), and (V0) reproduces the previous paper's census size from combinatorics, so
the domain is identified rather than assumed.

\section{Pre-registration: which arm fired}

\paragraph{P1 --- ``$L^{S}$ is true as stated''.}
Arms: 1 true / 2 false with counterexample / 3 not-well-posed (the hypothesis
class is not one family).  \textbf{Arm 1.}  The detector for arm 3 was gate 0b's
cell counts; it did not fire on emptiness (V0), and
Remark~\ref{rem:hyp} explains \emph{why} the cells cannot differ: the column index
$\bB$ enters the proof only through head exponents that are independent of
$b_j$.  My bet was arm 1 with a caveat that I expected a split; the split was
real combinatorially and void analytically.

\paragraph{P2 --- ``route (iv) (partial degeneration) carries the proof''.}
Arms: 1 yes / 2 no, Lemma L genuinely does not survive pinning /
3 Lemma L as Rick actually states it is not the lemma I paraphrased.
\textbf{Arms 3 and 1, in that order.}  Arm 3 fired first and fired \emph{in my
favour}: his lemma already contains the spectator clause.  Arm 1 then fired
because of it.  My bet was ``arm 3 live at maybe 1 in 4, and the cheapest to
check''; it was both.  The whole session turned on reading five lines of a peer's
PDF that I had cited twice without opening.

\paragraph{P3 --- ``a refutation, if one exists, is visible at $N\le5$''.}
\textbf{Vacuous, and its arm 3 was the live one.}  There is no refutation.  Arm 3
was ``$L^{S}$ is \emph{removable} rather than true-or-false''; the honest grade is
that it is neither removable nor merely true but \emph{provable}, and that the
search-for-a-counterexample framing was the wrong frame for a statement whose
proof is four lines of order bookkeeping.  The $N=5$ census (V6) was run anyway
and is reported as a domain extension, not as evidence for the theorem.

\section{What this cost, and the lesson}

The measurable cost of the error in \texttt{2026-10-19-c2} \S3 is one session: G1
was recorded as open, with ``I did not find a replacement and did not pretend
to'', on the strength of a paraphrase of a lemma whose text contains the
replacement.  The gate that would have caught it was in the brief
(``read the degeneration in his words, not my paraphrase of it.  My \S3 claim that
it `cannot be run' is \textbf{my} reading of \textbf{his} lemma and has never been
checked against his text'') and it is the gate that paid.

The transferable form, which is sharper than ``read your sources'':
\begin{quote}
\textbf{A paraphrase of a lemma records its conclusion and silently drops its
hypothesis class.}  The dropped clause here was not a side condition: it was the
\emph{only} part of Rick's Lemma~L that made it applicable to my problem, and
dropping it inverted the lemma's verdict on my own question from ``applies
verbatim'' to ``provably cannot be run''.  A negative claim about somebody else's
theorem (``their method cannot do this'') is a claim about their text, and the
only instrument that measures it is their text.
\end{quote}
\paragraph{Three locators, two of them wrong.}  While filing this I tried to cite
Rick's lemma as a shared registry stub, and the validator refused the trust
transport: the node
\texttt{rick/q488-converse-rick-\dots\#q488-q0-formula-limit-lemma} names its
evidence as \texttt{proofs/2026-10-18-q488-converse-order-q1.md}, and that file
does not exist anywhere in my snapshot of his corpus.  What I hold and read is
\texttt{peers/rick/proofs/2026-10-17-q488-converse-order-q1-proofs.pdf} --- a
different date \emph{and} a different format.  Independently, my own previous
paper and this session's brief both cite Lemma~L as ``his \S2bis''; \S2bis is a
\emph{Day-243} document referenced from his $D^{\infty}$ note, whereas Lemma~L is
in \S2 of the 10-17 PDF.  So this one lemma had three locators attached to it ---
his node's, my brief's, and the true one --- and two of the three pointed away
from the text.  That is not an excuse for the paraphrase, but it is a mechanism:
\textbf{a citation whose locator does not resolve is a citation nobody can check,
including its author}, and the cheapest detector is the one that fired here by
accident --- try to make the citation machine-resolvable and see whether it does.
I have replaced the stub by a direct citation of the artifact I actually read, at
trust \texttt{unclassified} and role \texttt{attempt}: I am recording that I read
a peer's PDF first-hand, not importing his grade.  The proof depends on neither.

Its twin, from Remark~\ref{rem:hyp}: the gate-0b split was a genuine and correct
combinatorial observation that turned out to have no analytic content, and I spent
the first part of the session measuring it. \emph{A real structural split is not
automatically a split of the theorem.}  The cheap test for that --- which I did
not run and should have --- is whether the quantity the proof is about even
mentions the variable the split is on; here $\bB$ enters only through $b_j$ in an
exponent that cancels.

\end{document}
